\ifdefined\gkver\else\def\gkver{student}\fi
\documentclass[\gkver]{gaokaozhenti}
\begin{document}

\chapter{2013年全国理综大纲卷}
%% paperType: 理综物理部分

\begin{enumerate}
\item
%% number: 14
%% paperName: 2013年全国理综大纲卷第 14 题 6 分
%% typeId: 1
%% type: 单选题
%% chapter: 光学
%% point: 光的衍射
%% method: 无
%% score: 6
%% degree: 850
%% duplicateId: 0
%% body:
下列现象中，属于光的衍射现象的是\xzanswer{B}

\fourchoices[answer=B]
{雨后天空出现彩虹}
{通过一个狭缝观察日光灯可看到彩色条纹}
{海市蜃楼现象}
{日光照射在肥皂泡上出现彩色条纹}

%% answer: B
%% memo:
\memoanswer{
本题原卷及材料中未提供详解，待补充。
}

\item
%% number: 15
%% paperName: 2013年全国理综大纲卷第 15 题 6 分
%% typeId: 2
%% type: 多选题
%% chapter: 功和能
%% point: 能量守恒定律
%% method: 无
%% score: 6
%% degree: 850
%% duplicateId: 0
%% body:
根据热力学定律，下列说法正确的是\xzanswer{AB}

\fourchoices[answer=AB]
{电冰箱的工作过程表明，热量可以从低温物体向高温物体传递}
{空调机在制冷过程中，从室内吸收的热量少于向室外放出的热量}
{科技的进步可以使内燃机成为单一热源的热机}
{对能源的过度消耗将使自然界得能量不断减少，形成能源危机}

%% answer: AB
%% memo:
\memoanswer{
本题原卷及材料中未提供详解，待补充。
}

\item
%% number: 16
%% paperName: 2013年全国理综大纲卷第 16 题 6 分
%% typeId: 1
%% type: 单选题
%% chapter: 原子和原子核
%% point: 半衰期
%% method: 无
%% score: 6
%% degree: 850
%% duplicateId: 0
%% body:
放射性元素氡（\ce{^{222}_{86}Rn}）经 $\alpha$ 衰变成为钋 \ce{^{218}_{84}Po}，半衰期为 3.8 天；但勘测表明，经过漫长的地质年代后，目前地壳中仍存在天然的含有放射性元素 \ce{^{222}_{86}Rn} 的矿石，其原因是\xzanswer{A}

\fourchoices[answer=A]
{目前地壳中的 \ce{^{222}_{86}Rn} 主要来自于其它放射元素的衰变}
{在地球形成的初期，地壳中元素 \ce{^{222}_{86}Rn} 的含量足够高}
{当衰变产物 \ce{^{218}_{84}Po} 积累到一定量以后，\ce{^{218}_{84}Po} 的增加会减慢 \ce{^{222}_{86}Rn} 的衰变进程}
{\ce{^{222}_{86}Rn} 主要存在于地球深处的矿石中，温度和压力改变了它的半衰期}

%% answer: A
%% memo:
\memoanswer{
本题原卷及材料中未提供详解，待补充。
}

\item
%% number: 17
%% paperName: 2013年全国理综大纲卷第 17 题 6 分
%% typeId: 1
%% type: 单选题
%% chapter: 电磁感应
%% point: 电磁感应中图象
%% method: 无
%% score: 6
%% degree: 650
%% duplicateId: 0
%% body:
纸面内两个半径均为 $R$ 的圆相切于 $O$ 点，两圆形区域内分别存在垂直纸面的匀强磁场，磁感应强度大小相等、方向相反，且不随时间变化。一长为 $2R$ 的导体杆 $OA$ 绕过 $O$ 点且垂直于纸面的轴顺时针匀速旋转，角速度为 $\omega$，$t=0$ 时，$OA$ 恰好位于两圆的公切线上，如图所示。若选取从 $O$ 指向 $A$ 的电动势为正，下列描述导体杆中感应电动势随时间变化的图像可能正确的是\xzanswer{C}

\onepicture[w=4.2cm, align=right]{figs/17a.png}

\fourchoices[answer=C, ispicture=true, h=2.6cm]
{figs/17b0.png}
{figs/17b1.png}
{figs/17b2.png}
{figs/17b3.png}

%% answer: C
%% memo:
\memoanswer{
本题原卷及材料中未提供详解，待补充。
}

\item
%% number: 18
%% paperName: 2013年全国理综大纲卷第 18 题 6 分
%% typeId: 1
%% type: 单选题
%% chapter: 曲线运动
%% point: 天体运动
%% method: 无
%% score: 6
%% degree: 650
%% duplicateId: 0
%% body:
“嫦娥一号”是我国首次发射的探月卫星，它在距月球表面高度为 $200\Ukm$ 的圆形轨道上运行，运行周期为 127 分钟。已知引力常量 $G=6.67\times10^{-11}\UNmqkgq$，月球的半径为 $1.74\times10^{3}\Ukm$。利用以上数据估算月球的质量约为\xzanswer{D}

\fourchoices[answer=D]
{$8.1\times10^{10}\Ukg$}
{$7.4\times10^{13}\Ukg$}
{$5.4\times10^{19}\Ukg$}
{$7.4\times10^{22}\Ukg$}

%% answer: D
%% memo:
\memoanswer{
本题原卷及材料中未提供详解，待补充。
}

\item
%% number: 19
%% paperName: 2013年全国理综大纲卷第 19 题 6 分
%% typeId: 2
%% type: 多选题
%% chapter: 直线运动
%% point: 竖直上下抛运动
%% method: 无
%% score: 6
%% degree: 650
%% duplicateId: 0
%% body:
将甲乙两小球先后以同样的速度在距地面不同高度处竖直向上抛出，抛出时间间隔为 2 s，它们运动的图像分别如直线甲乙所示。则\xzanswer{BD}

\onepicture[w=5.5cm, align=right]{figs/19.png}

\fourchoices[answer=BD]
{$t=2\Us$ 时，两球的高度相差一定为 $40\Um$}
{$t=4\Us$ 时，两球相对于各自的抛出点的位移相等}
{两球从抛出至落到地面所用的时间间隔相等}
{甲球从抛出至到达最高点的时间间隔与乙球相等}

%% answer: BD
%% memo:
\memoanswer{
本题原卷及材料中未提供详解，待补充。
}

\item
%% number: 20
%% paperName: 2013年全国理综大纲卷第 20 题 6 分
%% typeId: 2
%% type: 多选题
%% chapter: 功和能
%% point: 动能定理
%% method: 无
%% score: 6
%% degree: 650
%% duplicateId: 0
%% body:
如图所示，一固定斜面倾角为 $30^\circ$，一质量为 $m$ 的小物块自斜面底端以一定的初速度沿斜面向上做匀减速运动，加速度大小等于重力加速度的大小 $g$。物块上升的最大高度为 $H$，则此过程中，物块的\xzanswer{AC}

\onepicture[w=5.5cm, align=right]{figs/20.png}

\fourchoices[answer=AC]
{动能损失了 $2mgH$}
{动能损失了 $mgH$}
{机械能损失了 $mgH$}
{机械能损失了 $\frac{1}{2}mgH$}

%% answer: AC
%% memo:
\memoanswer{
本题原卷及材料中未提供详解，待补充。
}

\item
%% number: 21
%% paperName: 2013年全国理综大纲卷第 21 题 6 分
%% typeId: 1
%% type: 单选题
%% chapter: 机械波
%% point: 干涉衍射
%% method: 无
%% score: 6
%% degree: 450
%% duplicateId: 0
%% body:
在学校运动场上 $50\Um$ 直跑道的两端，分别安装了由同一信号发生器带动的两个相同的扬声器。两个扬声器连续发出波长为 $5\Um$ 的声波。一同学从该跑道的中点出发，向某一段点缓慢行进 $10\Um$。在此过程中，他听到扬声器声音由强变弱的次数为\xzanswer{B}

\fourchoices[answer=B]
{2}
{4}
{6}
{8}

%% answer: B
%% memo:
\memoanswer{
本题原卷及材料中未提供详解，待补充。
}

\item
%% number: 22
%% paperName: 2013年全国理综大纲卷第 22 题 6 分
%% typeId: 4
%% type: 实验题
%% chapter: 电场
%% point: 描绘等势面
%% method: 无
%% score: 6
%% degree: 650
%% duplicateId: 0
%% body:
如图所示，E 为直流电源，G 为灵敏电流计，A、B 为两个圆柱形电极，P 是木板，C、D 为两个探针，S 为开关。现用上述实验器材进行“用描迹法画出电场中平面上的等势线”的实验。

\begin{enumerate}
	\item
	木板 P 上有白纸、导电纸和复写纸，最上面的应该是\tkanswer[1.2cm]{导电}纸；
	\item
	用实线代表导线将实验器材正确连接。
\end{enumerate}
\onepicture[w=7.5cm, align=right]{figs/22.png}

%% answer:
\jdanswer{
\begin{enumerate}
	\item
	导电

	\item
	如图所示

	\onepicture[w=7.5cm]{figs/22_an.png}

\end{enumerate}
}
%% memo:
\memoanswer{
本题原卷及材料中未提供详解，待补充。
}

\item
%% number: 23
%% paperName: 2013年全国理综大纲卷第 23 题 12 分
%% typeId: 4
%% type: 实验题
%% chapter: 曲线运动
%% point: 研究平抛运动实验
%% method: 无
%% score: 12
%% degree: 650
%% duplicateId: 0
%% body:
测量小物块 Q 与平板 P 之间的动摩擦因数的实验装置如图所示。AB 是半径足够大的、光滑的四分之一圆弧轨道，与水平固定放置的 P 板的上表面 BC 在 B 点相切，C 点在水平地面的垂直投影为 $C'$。重力加速度为 $g$。实验步骤如下：

\begin{steps}
	\item
	用天平称出物块 Q 的质量 $m$；
	\item
	测量出轨道 AB 的半径 $R$、BC 的长度 $L$ 和 $CC'$ 的高度 $h$；
	\item
	将物块 Q 在 A 点由静止释放，在物块 Q 落地处标记其落地点 D；
	\item
	重复步骤③，共做 10 次；
	\item
	将 10 个落地点用一个尽量小的圆围住，用米尺测量圆心到 $C'$ 的距离 $s$。
\end{steps}

\begin{enumerate}
	\item
	用实验中的测量量表示：
	（i）物块 Q 到达 B 点时的动能 $E_{kB}=$\tkanswer[1.8cm]{$mgR$}；
	（ii）物块 Q 到达 C 点时的动能 $E_{kC}=$\tkanswer[2.6cm]{$\frac{mg s^2}{4h}$}；
	（iii）在物块 Q 从 B 运动到 C 的过程中，物块 Q 克服摩擦力做的功 $W_f=$\tkanswer[3.2cm]{$mgR-\frac{mg s^2}{4h}$}；
	（iv）物块 Q 与平板 P 之间的动摩擦因数 $\mu=$\tkanswer[3.0cm]{$\frac{R}{L}-\frac{s^2}{4hL}$}。

	\item
	回答下列问题：
	（i）实验步骤④⑤的目的是\tkanswer[2.6cm]{减小实验结果的误差}；
	（ii）已知实验测得的值比实际值偏大，其原因除了实验中测量量的误差外，其它问题可能是\tkanswer[3.0cm]{圆弧轨道存在摩擦（或接缝 B 处不平滑等）}。（写出一个可能的原因即可）
\end{enumerate}
\onepicture[w=5.5cm, align=right]{figs/23.png}

%% answer:
\jdanswer{
\begin{enumerate}
	\item
	（i）$mgR$；（ii）$\frac{mg s^2}{4h}$；（iii）$mgR-\frac{mg s^2}{4h}$；（iv）$\frac{R}{L}-\frac{s^2}{4hL}$（每空 2 分）

	\item
	（i）减小实验结果的误差（2 分）；（ii）圆弧轨道存在摩擦（或接缝 B 处不平滑等）（2 分）

\end{enumerate}
}
%% memo:
\memoanswer{
本题原卷及材料中未提供详解，待补充。
}

\item
%% number: 24
%% paperName: 2013年全国理综大纲卷第 24 题 15 分
%% typeId: 6
%% type: 计算题
%% chapter: 直线运动
%% point: 匀变速直线运动
%% method: 无
%% score: 15
%% degree: 850
%% duplicateId: 0
%% body:
一客运列车匀速行驶，其车轮在轨道间的接缝处会产生周期性的撞击。坐在该客车中的某旅客测得从第 1 次到第 16 次撞击声之间的时间间隔为 $10.0\Us$。在相邻的平行车道上有一列货车，当该旅客经过货车车尾时，火车恰好从静止开始以恒定加速度沿客车行进方向运动。该旅客在此后的 $20.0\Us$ 内，看到恰好有 30 节货车车厢被他连续超过。已知每根轨道的长度为 $25.0\Um$，每节货车车厢的长度为 $16.0\Um$，货车车厢间距忽略不计。求：

\begin{enumerate}
	\item
	客车运行的速度大小；
	\item
	货车运行加速度的大小。
\end{enumerate}

%% answer:
\jdanswer{
\begin{enumerate}
	\item
	$v=37.5\Ums$

	\item
	$a=1.35\Umsq$

\end{enumerate}
}
%% memo:
\memoanswer{
（1）设连续两次撞击轨道的时间间隔为 $\Delta t$，每根轨道的长度为 $l$，则客车的速度为 $v=\frac{l}{\Delta t}$（3 分），其中 $l=25.0\Um$，$\Delta t=\frac{10.0}{16-1}\Us$，解得 $v=37.5\Ums$（2 分）。

（2）设从货车开始运动后 $t=20.0\Us$ 内客车行驶的距离为 $s_1$，货车行驶的距离为 $s_2$，货车的加速度为 $a$，30 节货车车厢的总长度为 $L=30\times16.0\Um$。由运动学公式有 $s_1=vt$（3 分），$s_2=\frac{1}{2}at^2$（3 分），由题意有 $L=s_1-s_2$（2 分），联立解得 $a=1.35\Umsq$（2 分）。
}

\item
%% number: 25
%% paperName: 2013年全国理综大纲卷第 25 题 19 分
%% typeId: 6
%% type: 计算题
%% chapter: 电场
%% point: 带电粒子的往复运动
%% method: 无
%% score: 19
%% degree: 450
%% duplicateId: 0
%% body:
一电荷量为 $q$（$q>0$）、质量为 $m$ 的带电粒子在匀强电场的作用下，在 $t=0$ 时由静止开始运动，场强随时间变化的规律如图所示。不计重力，求在 $t=0$ 到 $t=T$ 的时间间隔内

\begin{enumerate}
	\item
	粒子位移的大小和方向；
	\item
	粒子沿初始电场反方向运动的时间。
\end{enumerate}
\onepicture[w=5.5cm, align=right]{figs/25.png}

%% answer:
\jdanswer{
\begin{enumerate}
	\item
	$\frac{qE_0T^2}{16m}$，方向沿初始电场正方向

	\item
	$\frac{T}{4}$

\end{enumerate}
}
%% memo:
\memoanswer{
解法一：粒子在 $0\sim\frac{T}{4}$、$\frac{T}{4}\sim\frac{T}{2}$、$\frac{T}{2}\sim\frac{3T}{4}$、$\frac{3T}{4}\sim T$ 时间间隔内做匀变速运动，设加速度分别为 $a_1$、$a_2$、$a_3$、$a_4$，由牛顿第二定律得 $qE_0=ma_1$，$2qE_0=-ma_2$，$2qE_0=ma_3$，$qE_0=-ma_4$……（每个式子 1 分）。

由此得带电粒子在 $0\sim T$ 时间间隔内运动的 $a-t$ 图像如图（a）所示（2 分），对应的 $v-t$ 图像如图（b）所示（3 分），

\twopicture[numA={（a）}, numB={（b）}, hA=3.6cm, hB=3.6cm, align=center]
{figs/25_an1.png}
{figs/25_an2.png}

其中 $v_1=a_1\frac{T}{4}=\frac{qE_0T}{4m}$（1 分）。由图（b）可知，带电粒子在 $t=0$ 到 $t=T$ 时的位移为 $s=\frac{T}{4}v_1$（2 分），联立解得 $s=\frac{qE_0T^2}{16m}$（2 分），它的方向沿初始电场正方向。（1 分）

（2）由图（b）可知，粒子在 $t=3T/8$ 到 $t=5T/8$ 内沿初始电场反方向运动，总的运动时间为 $t=\frac{5T}{8}-\frac{3T}{8}=\frac{T}{4}$（4 分）。

解法二：带电粒子在 $0\sim\frac{T}{4}$、$\frac{T}{4}\sim\frac{T}{2}$、$\frac{T}{2}\sim\frac{3T}{4}$、$\frac{3T}{4}\sim T$ 时间间隔内做匀变速运动，设加速度分别为 $a_1$、$a_2$、$a_3$、$a_4$，由牛顿第二定律得 $qE_0=ma_1$，$2qE_0=-ma_2$，$2qE_0=ma_3$，$qE_0=-ma_4$……（每个式子 1 分）。

设粒子在 $t=\frac{T}{4}$、$t=\frac{T}{2}$、$t=\frac{3T}{4}$、$t=T$ 时刻的速度分别为 $v_1$、$v_2$、$v_3$、$v_4$，则有 $v_1=a_1\frac{T}{4}$、$v_2=v_1+a_2\frac{T}{4}$、$v_3=v_2+a_3\frac{T}{4}$、$v_4=v_3+a_4\frac{T}{4}$（每个式子 1 分）。

设带电粒子在 $t=0$ 到 $t=T$ 时的位移为 $s$，有 $s=\left(\frac{v_1}{2}+\frac{v_1+v_2}{2}+\frac{v_2+v_3}{2}+\frac{v_3+v_4}{2}\right)\frac{T}{4}$（4 分），解得 $s=\frac{qE_0T^2}{16m}$（2 分），它的方向沿初始电场正方向。（1 分）

（2）由电场的变化规律知，粒子从 $t=\frac{T}{4}$ 时开始减速，设经过时间 $t_1$ 粒子速度为零，有 $0=v_1+a_2t_1$，解得 $t_1=\frac{T}{8}$（1 分）；粒子从 $t=\frac{T}{2}$ 时开始加速，设经过时间 $t_2$ 粒子速度为零，有 $0=v_2+a_3t_2$，解得 $t_2=\frac{T}{8}$（1 分）；设粒子从 $t=0$ 到 $t=T$ 内沿初始电场反方向运动的时间为 $t$，有 $t=\left(\frac{T}{4}-t_1\right)+t_2$（1 分），解得 $t=\frac{T}{4}$（1 分）。
}

\item
%% number: 26
%% paperName: 2013年全国理综大纲卷第 26 题 20 分
%% typeId: 6
%% type: 计算题
%% chapter: 磁场
%% point: 带电粒子在磁场中的运动
%% method: 无
%% score: 20
%% degree: 450
%% duplicateId: 0
%% body:
如图所示，虚线 $OL$ 与 $y$ 轴的夹角为 $\theta=60^\circ$，在此角范围内有垂直于 $xOy$ 平面向外的匀强磁场，磁感应强度大小为 $B$。一质量为 $m$、电荷量为 $q$（$q>0$）的粒子从左侧平行于 $x$ 轴射入磁场，入射点为 $M$。粒子在磁场中运动的轨道半径为 $R$。粒子离开磁场后的运动轨迹与 $x$ 轴交于 $P$ 点（图中未画出），且 $OD=R$。不计重力。求 $M$ 点到 $O$ 点的距离和粒子在磁场中运动的时间。
\onepicture[w=5.5cm, align=right]{figs/26.png}

%% answer:
\jdanswer{
$h=(1-\frac{\sqrt{3}}{3})R$，$t=\frac{\pi m}{6qB}$（$\alpha=30^\circ$）

$h=(1+\frac{\sqrt{3}}{3})R$，$t=\frac{\pi m}{2qB}$（$\alpha=90^\circ$）
}
%% memo:
\memoanswer{
根据题意，粒子进入磁场后做匀速圆周运动，设运动轨迹交虚线 $OL$ 于 $A$ 点，圆心在 $y$ 轴上的 $C$ 点，$AC$ 与 $y$ 轴的夹角为 $\alpha$；粒子从 $A$ 点射出后，运动轨迹交 $x$ 轴的 $P$ 点，设 $AP$ 与 $x$ 轴的夹角为 $\beta$，如图所示。有（判断出圆心在 $y$ 轴上得 1 分）

$qvB=m\frac{v^2}{R}$（1 分），周期为 $T=\frac{2\pi m}{qB}$（1 分）。

过 $A$ 点作 $x$、$y$ 轴的垂线，垂足分别为 $B$、$D$。由几何知识得 $\overline{AD}=R\sin\alpha$，$\overline{OD}=\overline{AD}\cot60^\circ$，$\overline{BP}=\overline{OD}\cot\beta$，$\overline{OP}=\overline{AD}+\overline{BP}$。

\onepicture[w=5.2cm, align=right]{figs/26_an.png}

$\alpha=\beta$（2 分），联立得到 $\sin\alpha+\frac{1}{\sqrt{3}}\cos\alpha=1$（2 分），解得 $\alpha=30^\circ$，或 $\alpha=90^\circ$（各 2 分）。

设 $M$ 点到 $O$ 点的距离为 $h$，有 $\overline{AD}=R\sin\alpha$，$h=R-\overline{OC}$，$\overline{OC}=\overline{CD}-\overline{OD}=R\cos\alpha-\frac{\sqrt{3}}{3}\overline{AD}$，联立得到 $h=R-\frac{2}{\sqrt{3}}R\cos(\alpha+30^\circ)$（1 分），解得 $h=(1-\frac{\sqrt{3}}{3})R$（$\alpha=30^\circ$）（2 分），$h=(1+\frac{\sqrt{3}}{3})R$（$\alpha=90^\circ$）（2 分）。

当 $\alpha=30^\circ$ 时，粒子在磁场中运动的时间为 $t=\frac{T}{12}=\frac{\pi m}{6qB}$（2 分）；当 $\alpha=90^\circ$ 时，粒子在磁场中运动的时间为 $t=\frac{T}{4}=\frac{\pi m}{2qB}$（2 分）。
}

\end{enumerate}

\end{document}
